\section{Exact transport of additional incidence counts}
\label{sec:passengers}

Complete no-fill schedules also simplify the transport of an additional
nonnegative incidence row. This is an algebraic refinement with a clear
use in later geometric work. It transports counts under the stated positive
update; it does not recover cyclic order, strand endpoints or an embedding.

Let $N$ now contain the \emph{exact uncapped} source counts, and let
$((r_0,g_0),\ldots,(r_{m-1},g_{m-1}))$ be a verified complete
column-leaf schedule. Let $p=(p_h)$ be an additional nonnegative row, called
a passenger. When the pivot $(r_i,g_i)$ is applied to it, the formal positive
transport rule is
\begin{equation}\label{eq:passenger-update}
 p'_h=p_h+p_{g_i}n_{r_i h}\quad(h\ne g_i),
 \qquad p'_{g_i}\text{ is deleted}.
\end{equation}
In a geometric application this formula must itself be justified by the
chosen disk slides and the placement of the passenger curves. The theorem
below is conditional on this exact update rule, not on abelianized words.

\begin{lemma}[Immutable donors form a directed acyclic graph]
\label{lem:donor-dag}
For $j<i$, the source entry $n_{r_i g_j}$ is zero. A donor row therefore
contains only its own pivot column, later pivot columns and survivor
columns. Apart from deleting its pivot entry, its exact source coefficients
can be used without any row update throughout the whole schedule.
\end{lemma}

\begin{proof}
When column $g_j$ was used, row $r_i$ was still live. The column-leaf
condition forces its entry in that row to be zero. Every preceding pivot
has no fill in any surviving donor, so the surviving row entries have
always agreed with the source. This proves both statements.
\end{proof}

Form a directed weighted graph on the matrix columns. For every nonpivot
entry of donor row $r_i$, put an edge $g_i\to h$ of weight $n_{r_i h}$.
Order the pivot columns by schedule time and put survivors afterwards.
Lemma~\ref{lem:donor-dag} makes this a directed acyclic graph, with survivor
columns as sinks. The pivot entries themselves do not give edges.

\Needspace{13\baselineskip}
\begin{theorem}[Sparse positive transport and a bit bound]
\label{thm:passenger}
One passenger row can be transported exactly by reading every nonpivot
donor incidence once, using $O(I+m+|\supp p|)$ big-integer operations.
For a survivor $h$, the final entry is
\begin{equation}\label{eq:passenger-paths}
 p_h^{\mathrm{out}}
 =\sum_{v}p_v\sum_{\gamma:v\leadsto h}
          \prod_{e\in\gamma}w_e,
\end{equation}
with the length-zero path included. Let $E$ be the number of donor graph
edges, let
\[
 L_N=\sum_{e}\left\lceil\log_2(1+w_e)\right\rceil,
 \qquad
 B_p=\max\left(\{0\}\cup
       \{\lceil\log_2(1+p_v)\rceil:v\in\supp p\}\right),
\]
and let $q$ be the number of appearing source or passenger columns.
Every intermediate and final nonzero passenger entry has bit length at
most
\begin{equation}\label{eq:passenger-bit}
 B_p+L_N+E+\lceil\log_2(q+1)\rceil+1.
\end{equation}
Hence transport is polynomial in the encoded exact source data, even when
the numerical coefficients grow exponentially in the number of rows.
\end{theorem}

\begin{proof}
At pivot $g_i$, distribute its current passenger coefficient along its
outgoing donor edges, multiplying by their weights, and then delete it.
The immutable-donor lemma shows that these are exactly the updates
\eqref{eq:passenger-update}. Every graph edge is processed once, giving
the structural operation count.

Induction in topological order identifies the value reaching a vertex
with the sum over all directed paths reaching that vertex. This proves
\eqref{eq:passenger-paths}; columns occurring only in the passenger have
only their length-zero contribution.

A directed path uses every edge at most once. Its edge set determines
the path, so there are at most $2^E$ paths from any fixed starting vertex
to a specified vertex. Each path product is at most $2^{L_N}$, and
each initial passenger value is less than $2^{B_p}$. Summing over at most
$q$ starting vertices gives the stated conservative bit bound. An
intermediate entry is a sub-sum of the analogous nonnegative path
expression ending at that vertex, so it satisfies the bound as well.
Polynomial-time integer multiplication
then gives polynomial bit complexity.
\end{proof}

\begin{example}[Two branching paths]
Use columns $(g_0,g_1,g_2,z)$ and the source rows
\[
 N=\begin{pmatrix}
 1&2&3&0\\
 0&1&0&5\\
 0&0&1&7
 \end{pmatrix},\qquad
 p=(11,0,0,13).
\]
The complete column schedule uses $g_0,g_1,g_2$. The final survivor entry is
\[
 p_z^{\mathrm{out}}=13+11(2\cdot5+3\cdot7)=354.
\]
Using capped donor counts would instead give $101$. Capping is exact for
the cancellation \emph{decision}, but not for the transported values.
\end{example}

The delivered \code{raw_passenger_transport} module first checks the full
schedule using the independent source timestamp verifier, then applies
the uncapped donor updates. Its tests compare with an independent explicit
positive-pivot oracle. The API includes columns that appear only in the
passenger and retains them as survivors.

This gives a concrete polynomial-bit algebraic step toward tracking
additional curves through a successful cancellation. It is intentionally
short of geometric boundary-pattern transport. Two patterns with identical
unsigned incidence can have different cyclic orders and topology.
Moreover, the virtual band columns of Section~\ref{sec:saturated-bundles}
were omitted for a decision proof; this theorem does not reconstruct
per-band passenger data from a genus counter. A compressed geometric
representation of those additional markings remains a research problem.
